% Raw pre-archive reasoning, 22:18 UTC. Not promoted or claimed proved.
\subsection*{Candidate: three cubic components instead of principal theta}
Let $B$ be an indecomposable nonordinary ppav surface, and let
$\phi:B\to A$ be any odd polarized isogeny of multiplier $n$,
with kernel $K$ of order $k=n^2$. Use the already proved full
source monodromy on $|9\Theta_B|$, of relative rank $80$.
The proposed norm target has odd relative rank $81k-1$.
No assumption that the individual principal theta norm is nodal
should be necessary if the following argument is completed.

Take three generic smooth cubic divisors $C_i\in|3\Theta|$.
Each individual norm is generically nodal by the proved odd-natural
block incidence, for arbitrary $K$. Choose the three so all shifted
cross intersections and self-nodes are clean. A selected simple
tangency between $C_i$ and $C_j$ can be imposed while retaining the
other norm identifications as distinct ordinary nodes.

The contact incidence for $C_i(P)$ and $C_j(P-h)$ with common
tangent direction is a $\mathbf P^6\times\mathbf P^6$ bundle over
the projective tangent bundle of $B$, hence irreducible of dimension
$15$. Both projections to the cubic coefficient spaces are
dominant, so smoothness and individual norm-nodality are dense opens.
The single-point second-jet theorem gives exact ordinary contact.
For two contacts at $P,Q$, generic joint first jets give eight
independent equations against six point/direction parameters.
If either joint-jet pair lies in the exceptional set $R$, the base
point-pair locus has dimension at most two. After the two values,
the gradient space has dimension at least two, with nonzero
projection to each tangent space. The locus of direction pairs
on which the two derivative rows have rank $s$ has dimension at
most $s$ for $s=0,1,2$. Thus even if both joint-jet pairs are bad,
the double-contact incidence has dimension at most
$16-4-s_1-s_2+2+\min(s_1,s_2)\le14$.
Coincident contact points force the two directions equal on a
smooth component, reducing the same count. This excludes all
unwanted shifted tangencies.

For a triple norm packet split between two components, the generic
count is seven value/directional conditions against five parameters.
The component with two extra values uses three-value independence.
If its directional derivative becomes dependent, either there is
only one bad direction (saving one parameter) or its tangent plane
is the span of the three value vectors. The latter puts one
point-pair in $R$, saving two point parameters. A simultaneous
loss on the other component is again compensated by that same
two-dimensional exceptional locus. Hence the locus has dimension
at most fourteen. Packets meeting the marked contact image are
avoided by the finitely many extra values at its $K$-translates.
Choose the third cubic generically afterward: the finite singular
and pair-intersection sets are avoided, and its two-point norm
packets meet either fixed curve in a proper incidence.
These counts need full written audit before use.

On $|9\Theta|$, any three first jets are independent: for a chosen
point, square a cubic section vanishing at the other two and
nonzero there, then multiply by a cubic first-jet section. At four
points independence of cubic values gives the same construction;
for two fixed translated pairs its exceptions are finite by the
proved four-value lemma. The node covers have degree $162$ per
unordered kernel type and irreducible one-/two-point incidences.
The above clean boundary gives a single transposition in each
type and no others, yielding the full product $S_{162}^{(k-1)/2}$.
The source $\mathrm{GL}_{80}(\mathbf F_2)$ has no composition factor
in that product ($v_2|\mathrm{GL}_{80}|=3160$ whereas
$v_2|S_{162}|=159$), so Goursat and the proved natural-module
$H^1$ calculation give an actual source block.

At a clean three-cubic boundary the graph has three vertices,
eighteen loops per vertex and nonzero unordered type, and eighteen
cross edges for each vertex pair and each $g\in K$. Thus
\[
 \dim G=81k-29,\qquad
 \dim T_r=81(k-1),\qquad \dim T_{r,h}=162.
\]
The selected graph has $54$ edges and rank $52$.
The generic source relative rank is $80$, but its boundary rank is
$76+2f$; this rank drop is handled only by the proved clopen Lang
closure and finite torsor extension, never by a full-fibre isomorphism.
The earlier three-theta graph argument works with eighteen in place
of two: all type hearts span the codimension-one total-edge-sum
kernel $G^0$, and an $S_{162}^{(k-1)/2}$-invariant subspace contained
in every partial-normalization kernel $T_r$ is zero.

The key gluing refinement is linear, not a new rank-preservation
claim. At a selected ordinary tangency the inertia difference has
pure graph image $q=e_a+e_b$ in both source and target. The
equivariant splitting of the quotient Lang torsor therefore sends
$q$ to the same pure pair vector modulo its translation submodule.
The selected $S_{18}$ conjugates give every pair difference;
their sums include the total eighteen-edge vector. The two selected
pair totals at a vertex add to zero in source graph cohomology.
Their target lifts add to the unselected vertex star. In every
type this star has $72$ ones and $90$ zeros, so the independent
$S_{162}$ action forces the whole type heart into the translation
module. This resolves the compatibility issue provided all those
isolated tangencies belong to the same normal node-family base.

The usual zero/surjective projection dichotomy, clopen graph
comparison and absence of odd invariant covectors then give
irreducibility. The block plus all hearts has dimension
$80+161(k-1)/2>(81k-1)/2$, so the normal transvection subgroup is
irreducible; the actual natural block excludes an alternating form.
This would prove full $\mathrm{GL}_{81k-1}$ in target degree $3n$.

For every fixed two-rank-zero surface, the explicit geometric Q8
and the already proved four-squares descent supply self-isogenies
of every odd multiplier. The proposed transfer would therefore
give every odd degree divisible by three. Together with the known
odd degrees one and five, even addition would cover every even
degree: residue zero modulo six uses $3+3n$, residue four uses
$1+3n$, residue two uses $5+3n$, with degree two treated separately.
This is still a candidate assembly, not a promoted consequence.
